\section*{第 \ref{ch:b1:rays-reflection-refraction} 章——几何光学：光线、反射、折射}

\begin{solution}{exo:b1:rays-reflection-refraction:1}
$v = c/n$：水 \qty{2.26e8}{m/s}，玻璃 \qty{2.00e8}{m/s}，金刚石
\qty{1.24e8}{m/s}。光纤：$t = nL/c = 1.48 \times 10^3/3.00\times10^8 =
\qty{4.9}{\micro s}$.
\end{solution}

\begin{solution}{exo:b1:rays-reflection-refraction:2}
$\sin i_2 = \sin i_1/1.33$: $30^\circ \to 22.1^\circ$; $60^\circ \to
40.6^\circ$。从水中以 $50^\circ$ 入射：$1.33\sin 50^\circ = 1.02 > 1$，
没有折射光线——发生全反射（$i_c = 48.8^\circ$）。
\end{solution}

\begin{solution}{exo:b1:rays-reflection-refraction:3}
$i_c = \arcsin(1/n)$：玻璃 $41.8^\circ$，金刚石 $24.4^\circ$，水
$48.8^\circ$。天空圆盘的张角为 $2 \times 48.8^\circ = 97.5^\circ$。
\end{solution}

\begin{solution}{exo:b1:rays-reflection-refraction:4}
眼睛位于 \qty{1.60}{m} 高处。从脚部发出的光线
于眼高一半处（即 \qty{0.80}{m}）反射后到达眼睛；从头顶发出的光线
在头顶与眼睛的中点处反射，反射点高度为 \qty{1.65}{m}。
镜子从 \qty{0.80}{m} 延伸到 \qty{1.65}{m}：高度为 \qty{0.85}{m}，恰为
身高的一半——与距离无关（中点规则
在任何距离下都成立）。
\end{solution}

\begin{solution}{exo:b1:rays-reflection-refraction:5}
小角度近似下，从光疏介质中观察时视位置为 $h' = h/n$，
从光密介质中观察时为 $h' = nh$。硬币：$1.2/1.33 =
\qty{0.90}{m}$。渔夫：$1.33 \times 1.8 = \qty{2.4}{m}$。理由如下：
对于界面上水平距离为 $x$ 的点，$\tan i_2 \approx i_2 =
x/h$，$\tan i_1 \approx i_1 = x/h'$，且 $i_1 = n i_2$。
\end{solution}

\begin{solution}{exo:b1:rays-reflection-refraction:6}
$\sin\frac{A + D_{\min}}{2} = 1.52\sin 30^\circ = 0.760$，于是
$\frac{A + D_{\min}}{2} = 49.5^\circ$, $D_{\min} = 38.9^\circ$.
出射需要 $r' \leq i_c$；又由 $\sin r = \sin i/n \leq 1/n$ 可得
$r \leq i_c$，因此 $A = r + r' \leq 2i_c$。此处 $i_c = \arcsin(1/1.52) =
41.1^\circ$，$2i_c = 82.3^\circ > 60^\circ$：成立。
\end{solution}

\begin{solution}{exo:b1:rays-reflection-refraction:7}
$\mathrm{NA} = \sqrt{1.50^2 - 1.48^2} = \sqrt{0.0596} = 0.244$,
$\theta_a = 14.1^\circ$. $\Delta t/L = (1.50/c)(1.50/1.48 - 1) =
5.0\times10^{-9} \times 0.0135 = \qty{68}{ns/km}$。在 \qty{1}{km} 范围内：
$B_{\max} \approx 1/(2 \times \qty{68}{ns}) = \qty{7.4}{Mbit/s}$.
\end{solution}

\begin{solution}{exo:b1:rays-reflection-refraction:8}
$D \approx (n - 1)A$：红光 $0.510 \times 10^\circ = 5.10^\circ$，蓝光
$5.30^\circ$；光谱宽度为 $0.20^\circ$。
\end{solution}

\begin{solution}{exo:b1:rays-reflection-refraction:9}
进入时：$\sin r = \sin 45^\circ/1.5 = 0.471$，$r = 28.1^\circ$。两个
界面平行，因此在第二个界面上的入射角为 $r$，且
$\sin i' = n\sin r = \sin i$：$i' = i$，出射光线
与入射光线平行。沿长为 $e/\cos r$ 的内部路径，
垂直方向的偏移为 $(e/\cos r)\sin(i - r)$：
$d = 2.0 \times \sin 16.9^\circ/\cos 28.1^\circ = \qty{0.66}{cm}$.
\end{solution}

\begin{solution}{exo:b1:rays-reflection-refraction:10}
两次折射（$2(i - r)$）和两次反射（$2(\pi - 2r)$）：
$D = 2\pi + 2i - 6r$。$\dd D/\dd i = 2 - 6\cos i/(n\cos r) = 0$ 给出
$3\cos i = n\cos r$, $9\cos^2 i = n^2 - 1 + \cos^2 i$, $\cos^2 i =
(n^2 - 1)/8$. $n = 1.333$: $\cos^2 i = 0.0971$, $i = 71.8^\circ$,
$\sin r = 0.950/1.333$, $r = 45.5^\circ$, $D = 360^\circ + 143.7^\circ -
272.8^\circ = 230.9^\circ$：虹位于距反日点 $230.9^\circ - 180^\circ =
51^\circ$ 处，在主虹（$42^\circ$）之外。
红光（$n = 1.331$）位于 $50.4^\circ$，紫光（$1.343$）位于 $53.6^\circ$：
颜色顺序颠倒，紫光在外侧。
\end{solution}

\begin{solution}{exo:b1:rays-reflection-refraction:11}
$\sin i_c = 1.00026/1.00029 = 0.99997$，$i_c = 89.56^\circ$：光线
与路面的夹角在 $0.44^\circ = \qty{7.7}{mrad}$ 以内的光线会被全
反射（一般地，掠射角 $\approx \sqrt{2\Delta n/n}$）。从
\qty{1.2}{m} 高处看去，视线在超过
$1.2/0.0077 \approx \qty{160}{m}$ 的距离处以此角度掠射：那里的路面反射天空——
即“水洼”。
\end{solution}

\begin{solution}{exo:b1:rays-reflection-refraction:12}
$t(x) = \dfrac{\sqrt{a^2 + x^2}}{v_1} + \dfrac{\sqrt{b^2 + (d - x)^2}}{v_2}$;
$t'(x) = \dfrac{x}{v_1\sqrt{a^2 + x^2}} - \dfrac{d - x}{v_2\sqrt{b^2 + (d - x)^2}}
= \dfrac{\sin i_1}{v_1} - \dfrac{\sin i_2}{v_2}$ 为零，当且仅当
$\sin i_1/v_1 = \sin i_2/v_2$，即 $n_1\sin i_1 = n_2\sin i_2$，其中
$n = c/v$.
\end{solution}

\begin{solution}{pb:b1:rays-reflection-refraction:1}
\textbf{1.} $n = c/v$; $v_1 = 3.00\times10^8/1.480 = \qty{2.03e8}{m/s}$.

\textbf{2.} $t = Ln_1/c = 5.0\times10^4 \times 1.48/3.00\times10^8 =
\qty{247}{\micro s}$.

\textbf{3.} $\nu = c/\lambda_0 = \qty{1.94e14}{Hz}$，在玻璃中保持不变；
$\lambda = \lambda_0/n_1 = \qty{1047}{nm}$.

\textbf{4.} 裸露的纤芯只有在表面
完全干净、未被触碰时才能导光：灰尘、水或手握处会局部以更高折射率的物质取代空气，
使光线泄漏；划痕则会散射光线。包层
是一个密封、可控的界面，其折射率阶跃是经过刻意选择的。

\textbf{5.} $\sin i_c = 1.465/1.480 = 0.9899$, $i_c = 81.8^\circ$.

\textbf{6.} $90^\circ - 81.8^\circ = 8.2^\circ$.

\textbf{7.} $\sin\theta = n_1\sin\theta'$；导光当且仅当 $\cos\theta' \geq
n_2/n_1$，即 $\sin\theta' \leq \sqrt{1 - n_2^2/n_1^2}$，即
$\sin\theta \leq \sqrt{n_1^2 - n_2^2} = \sqrt{2.1904 - 2.1462} = 0.210$：
$\theta_a = 12.1^\circ$。

\textbf{8.} $\mathrm{NA} = 0.210$。阶跃越大，锥角越宽，
因此发散光源（发光二极管）所发出的光中被捕获的部分越多。

\textbf{9.} 光线与半径为 $R_b$ 的圆（内缘）相切，
并与半径为 $R_b + 2a$ 的圆（外缘）相交；在直角三角形“圆心—切点—入射点”中，
入射点处的
光线与半径（法线）之间的夹角满足
$\sin i = R_b/(R_b + 2a)$。导光当且仅当 $\sin i \geq n_2/n_1$：
即 $R_b \geq 2a\,\dfrac{n_2/n_1}{1 - n_2/n_1} = \qty{50}{\micro m} \times 97.6
= \qty{4.9}{mm}$。弯曲半径更小则会泄漏。

\textbf{10.} $t_0 = Ln_1/c$；最陡光线的路径长为 $L/\cos\theta'_{\max} =
Ln_1/n_2$，因此 $t_1 = Ln_1^2/(n_2 c)$。

\textbf{11.} $\Delta t/L = (n_1/c)(n_1/n_2 - 1) = 4.93\times10^{-9} \times
0.01024 = \qty{50.5}{ns/km}$；传输 \qty{2}{km} 时：\qty{101}{ns}。

\textbf{12.} $B_{\max} \approx 1/(2\Delta t) = 1/\qty{202}{ns} \approx
\qty{5}{Mbit/s}$。

\textbf{13.} $B_{\max} L = 1/(2\,\Delta t/L) = 1/(2 \times
\qty{50.5}{ns/km}) \approx \qty{10}{Mbit.s^{-1}.km}$。

\textbf{14.} 远离轴线的光线路径更长，但所处介质的折射率较低，
因此速度更快；若选取恰当的折射率分布，两种效应相互抵消，
所有光线同时到达。

\textbf{15.} $\qty{9}{\micro m} \approx 6\lambda_0$：纤芯相对于波长并不大，
因此衍射占主导，光线的概念已失去意义。
波动理论给出单一导模的图样：完全不存在模间展宽，
仅余色散与
衰减仍然存在。

\textbf{16.} \qty{100}{km} $\times$ \qty{0.2}{dB/km} $= \qty{20}{dB}$：
$P/P_0 = 10^{-2}$，即百分之一。

\textbf{17.} $10\log(P_0/P_{\min}) = \qty{30}{dB}$，因此 $L_{\max} =
30/0.20 = \qty{150}{km}$。

\textbf{18.} \qty{100}{km} 损耗 \qty{20}{dB}：余量为 \qty{10}{dB}，
即功率相差 $10$ 倍。

\textbf{19.} 发光二极管：$\Delta t_c = 17 \times 40 \times 100 = \qty{68000}{ps}
= \qty{68}{ns}$，$B \approx 1/\qty{136}{ns} = \qty{7}{Mbit/s}$。激光器：
$\Delta t_c = \qty{170}{ps}$，$B \approx 1/\qty{340}{ps} =
\qty{2.9}{Gbit/s}$。

\textbf{20.} $30/2.0 = \qty{15}{km}$，缩短为原来的十分之一。二氧化硅
在 \qty{1550}{nm} 附近最为透明（衰减最小），
且该处存在光放大器。

\textbf{21.} $30/30 = \qty{1}{km}$：放大器之间的间距为 \qty{1}{km}，
与 \qty{150}{km} 相比，光纤的优势为 $150$ 倍。

\textbf{22.} 模间展宽：\qty{101}{ns}（$\to$ \qty{5}{Mbit/s}）；
色散（激光器）：$17 \times 0.1 \times 2 = \qty{3.4}{ps}$；
衰减 \qty{0.4}{dB}，可忽略。限制速率的是模间展宽：
\qty{5}{Mbit/s}。

\textbf{23.} 不存在模间展宽；色散 \qty{170}{ps} $\to$
\qty{2.9}{Gbit/s}；衰减 \qty{20}{dB}，在预算之内。
限制速率的是色散，约为 \qty{3}{Gbit/s}。

\textbf{24.} $80 \times 2.9 \approx \qty{230}{Gbit/s}$。

\textbf{25.} 阶跃折射率：\qty{10}{Mbit.s^{-1}.km}；单模信道：
$2.9 \times 100 = \qty{290}{Gbit.s^{-1}.km}$——约为
$3\times10^4$ 倍。要点：只有一条导光路径，
因此不存在几何（模间）引起的到达时间展宽。
\end{solution}
